\section{Generátor vnitřního kódu (TAC)}

Modul \texttt{tac}(Three Address Code) realizuje transformaci abstraktního syntaktického stromu (AST) do tříadresového kódu (TAC). Tento mezikód slouží jako přechodná reprezentace mezi syntaktickou/semantickou analýzou a generováním cílového kódu (IFJcode25). Tříadresný kód používáme proto, že nám zjednoduší překlad do cílového kódu a umožní lepší optimalizaci. Generátor je implementován ve souborů \texttt{tac.c}.


\subsection*{Proces generování TAC}

Vstupním bodem je funkce \texttt{generate\_tac(...)}. Tato funkce přijímá kořen AST a rekurzivně prochází jeho uzly. Pro každý typ uzlu (např. přiřazení, podmínky, smyčky, volání funkcí) generuje odpovídající sadu TAC instrukcí.

Během průchodu stromem modul provádí následující klíčové operace:
\begin{enumerate}
    \item \textbf{Linearizace výrazů:} Rozklad složených výrazů na elementární instrukce s dočasnými proměnnými(např. \texttt{\$t1}, \texttt{\$t2}).
    \item \textbf{Řízení toku:} \texttt{if-else} a cyklus \texttt{while} jsou převedeny na lineární kód pomocí instrukcí skoků (\texttt{OP\_JUMP}, \texttt{OP\_JUMP\_IF\_FALSE}) a návěští (\texttt{OP\_LABEL}). Funkce \texttt{create\_unique\_label} generuje unikátní názvy návěští (např. \texttt{L\_ELSE\_0}, \texttt{L\_END\_IF\_1}).
    \item \textbf{Funkce a volání:} Pro definice funkcí generuje speciální instrukce \texttt{OP\_FUNCTION\_BEGIN} a \texttt{OP\_FUNCTION\_END}.Parametry funkcí jsou zpracovány a předány jako argumenty instrukcí.
\end{enumerate}

\subsection*{Použité datové struktury}

Výstupem generátoru je dvousměrně vázaný seznam (\texttt{TACDLList}). Tato struktura reprezentuje seznam pro tzv. čtveřici (quadruple):
\begin{itemize}
    \item \texttt{operation\_code}: Typ operace.
    \item \texttt{result}, \texttt{arg1}, \texttt{arg2}: Operandy instrukce.
\end{itemize}

Operandy jsou reprezentovány strukturou \texttt{Operand}, která může nést různé typy dat definované výčtem \texttt{OperandType}. 
Díky přímému propojení operandů typu \texttt{SYMBOL} se záznamy v tabulce symbolů (\texttt{TableEntry}) nemusí generátor cílového kódu provádět opakované vyhledávání názvů, což zvyšuje efektivitu překladu.